\section{\label{ch:udf}User-defined functions in constraints}

An arithmetic constraint expression\index{arithmetic constraint expression} is
normally built from integers, domain variables, and a fixed set of arithmetic
function symbols: \texttt{+}, \texttt{-}, \texttt{*}, \texttt{/},
\texttt{//}, \texttt{div}, \texttt{mod}, \texttt{**}, \texttt{abs},
\texttt{min}, \texttt{max}, \texttt{sum}, \texttt{prod},
\texttt{count}, and \texttt{cond}.  (In addition, index notations,
array comprehensions, and list comprehensions are interpreted as function
calls.)  Any other call that appears in a constraint expression is left as an
unevaluated\index{unevaluated} term, so a user-written function such as
\texttt{sqr(X) = X*X} cannot be used directly in a constraint.

The \texttt{udf}\index{\texttt{udf}} module (\texttt{lib/udf.pi}) adds
support for user-defined functions in constraints.  It is a
self-contained, extensible layer that works with all four constraint
solvers (\texttt{cp}, \texttt{sat}, \texttt{smt}, and \texttt{mip}).
A function is \emph{registered} in a symbolic registry, used as a term in
a constraint, and \emph{expanded} into primitive arithmetic before the
constraint is handed to the solver.

\subsubsection{The dollar-sign requirement}

The \texttt{udf} functions are not Picat functions: they exist only in the
registry and are matched as \emph{terms}\index{term}.  In a normal
(non-constraint) expression, \texttt{f(X)} is an ordinary call to a
function \texttt{f/1}, which does not exist.  Consequently, every use
--- in \texttt{define} heads, in \texttt{define} bodies, and at the call
sites of \texttt{flatten}, \texttt{evaluate}, and \texttt{post} --- must
be written with the term constructor \texttt{\$}.  Only the constraint
operators (\texttt{\#=}, etc.) treat expressions as terms, and re-implementing
those operators is precisely what the \texttt{udf.ops.pi} shims do.

\subsubsection{The built-ins of the \texttt{udf} module}

\begin{itemize}

\item \texttt{define($Head$,$Body$)}: Register a user-defined function.
      $Head$ is a pattern of the form \texttt{\$f($Arg_1$,\ldots,$Arg_n$)},
      and $Body$ is the function body given as a \emph{term} (prefix
      \texttt{\$}).  A function may have several (pattern) rules, and a body
      may call other \texttt{udf} functions or itself (recursion), up to a
      configurable expansion bound.  Example:
      \texttt{udf.define(\$total([]), 0)},
      \texttt{udf.define(\$total([H|T]), \$(H + total(T)))}.

\item \texttt{flatten($E$) = $Out$}: Rewrite the constraint expression $E$,
      expanding every registered user-function call into primitive
      arithmetic.  The result $Out$ uses only the primitive function symbols,
      so it can be used in any of the four solvers.

\item \texttt{evaluate($E$) = $Val$}: Like \texttt{flatten/1}, but also
      evaluates the result to a number.  This requires the expanded
      expression to be \emph{ground} (free of domain variables), so
      \texttt{evaluate} is only meaningful outside the constraint world:
      \texttt{udf.evaluate(\$total([1,2,3]))} is \texttt{6}.  For
      domain variables, use \texttt{\#=}, \texttt{post}, or the operator
      shims instead.

\item \texttt{post($Rel$,$E_1$,$E_2$)}: Expand $E_1$ and $E_2$ and
      post the constraint $E_1$ $Rel$ $E_2$, where $Rel$ is one of
      \verb+#=+, \verb+#!=+, \verb+#<+, \verb+#>+, \verb+#=<+,
      \verb+#>=+ (and the synonym \verb+#<=+).
      \texttt{udf.post('\#'=', Y, \$f(X))} is the explicit,
      zero-boilerplate style.

\item \texttt{set\_max\_depth($N$)}: Bound on how deep a (possibly
      recursive) body may expand before the exception
      \texttt{udf\_expansion\_too\_deep}\index{exception} is thrown.
      The default is \texttt{100}.

\end{itemize}

\subsubsection{The normal operator style with \texttt{udf.ops.pi}}

A solver operator cannot be overridden by an imported module, but a
module's \emph{own} definitions take precedence over an imported solver's.
which solver.  The include files
\texttt{lib/udf.ops.pi} (cp), \texttt{udf.ops\_sat.pi},
\texttt{udf.ops\_mip.pi}, and \texttt{udf.ops\_smt.pi}
each provide a one-line shim for every arithmetic and Boolean constraint
operator of that solver
(\verb+#= #!= #< #=< #<= #> #>=+ and
\verb+#~ #/\\ #^ #\\/ #=> #<=>+).  Adding
(\texttt{include "udf.ops.pi"}) to the \emph{user's own} module makes
the normal source style work:

\begin{verbatim}
    import cp.
    import udf.
    include "udf.ops.pi".

    main =>
        udf.define($day_demand(N), $(N*2+1)),
        P1 :: 0..20,
        P1 #>= $day_demand(1),          %  P1 >= 3
        ...
\end{verbatim}

\subsubsection{Examples and how to run them}

Two examples are provided in each of \texttt{exs/cp/},
\texttt{exs/sat/}, \texttt{exs/mip/}, and \texttt{exs/smt/}
(\texttt{udf\_functions.pi}); \texttt{exs/cp/udf\_ops\_schedule.pi}
adds a three-day production/stock scheduling problem that minimises
total production, using a recursive user function
\texttt{day\_demand/1} under \texttt{\#=} and \texttt{\#>=}.
\texttt{PICATPATH}\index{PICATPATH}.  From the \texttt{exs/cp} directory:

\begin{verbatim}
    picat -path <picat>/lib  udf_functions.pi
    PICATPATH=<picat>/lib  <picat>/emu/picat  udf_ops_schedule.pi
\end{verbatim}

where \texttt{<picat>/lib} is the repository's \texttt{lib} directory
(\texttt{udf.ops.pi}, which the examples include, sits next to them and is
found directly).

\subsubsection{Limitations}

Registering a function with \texttt{udf.define} makes it usable in
constraints, but \texttt{udf.evaluate} requires the expanded expression to
be ground.  Recursion is expanded structurally and is bounded by
\texttt{set\_max\_depth}; a body that does not structurally shrink (for
example, arithmetic recursion such as \texttt{f(N) = f(N-1)+1}) throws
\texttt{udf\_expansion\_too\_deep}.  For such cases, write the function to
reduce structurally (for instance, over a list).
